\par\noindent\textbf{(1)} 相位 \(\Phi(k)=kx-\omega(k)t\)。驻相点满足 \(\Phi'(k_*)=x-t\omega'(k_*)=0\)，故 \(x/t=v_g(k_*)\)。因此远场主贡献由群速等于观察射线速度的谱分量形成，并非一定取相速度 \(\omega/k\)。
\par\noindent\textbf{(2)} 若 \(\omega''(k_*)\ne0\)，二阶展开 \(\Phi\approx\Phi_*-\frac12t\omega''(k_*)(k-k_*)^2\)。Gaussian振荡积分得 \(\psi\sim A(k_*)\sqrt{2\pi/[t|\omega''(k_*)|]}\exp[i\Phi_*-i\frac\pi4\operatorname{sgn}\omega''(k_*)]\)，振幅通常按 \(t^{-1/2}\) 衰减。
\par\noindent\textbf{(3)} 若 \(\omega''(k_*)=0\) 而 \(\omega'''(k_*)\ne0\)，二阶驻相不再适用，须保留三阶相位 \(-t\omega'''(k_*)(k-k_*)^3/6\)。重缩放 \(\Delta k\sim t^{-1/3}\) 引出Airy型统一近似与 \(t^{-1/3}\) 的不同振幅尺度。
\par\noindent\textbf{(4)} 若 \(\omega(k)=ck\)，则 \(\psi(x,t)=\int A(k)e^{ik(x-ct)}dk=\psi(x-ct,0)\)，形状不变。这时 \(\omega''=0\) 在所有 \(k\) 上成立，但也没有非零三阶项，不能把普通驻相 \(t^{-1/2}\) 或Airy \(t^{-1/3}\) 套用到纯平移波。
